\subsection{Security Best Practices}\label{sec:security_of_cryptographic_algorithms:best_practices}

We now introduce the main high-level principles and methods for implementing cryptographic algorithms securely by mitigating and minimizing their exposure to the attacks described in \Cref{sec:security_of_cryptographic_algorithms:attacks}. Intuitively, the adoption of the following principles and methods cannot be considered sufficient to prevent all attacks.

\remarkBlock{Security best practices and programming languages}{The choice of the programming language for implementing cryptographic algorithms has a great impact on the relevance, practicality, and impact of the security best practices presented below --- see \Cref{sec:programming_languages_and_software:languages}.}

\subsubsection{Secret-Independent Execution}\label{sec:security_of_cryptographic_algorithms:best_practices:independent}

Ideally, no side-channel attack is possible if every observable feature of the execution of sensitive cryptographic computations (e.g., mathematical operations with variable complexity, branching, access to data) is \textbf{statistically independent of secret data}. In particular, \textbf{variable complexity} is a prominent cause of timing and power-related attacks, as it provides a direct correlation between observable behavior and secret data. Similarly, \textbf{branching} can reveal information either indirectly due to instruction caching --- which may cause timing leakage --- or directly by varying the amount or type of operations performed --- which may cause any kind of leakage. Finally, \textbf{data access} depending on secret data (e.g., as discussed for lookup tables in \gls{aes} --- see \Cref{sec:security_of_cryptographic_algorithms:attacks:sca}) can cause timing and cache leakages due to the (unpredictable) physical management of virtual memory. For instance, large arrays --- whose elements are contiguous in virtual memory --- could be actually stored on multiple pages that are mapped to different physical memory (and even in different chips) and consequently have different access times.

\remarkBlock{On the difficulty of achieving secret-independent execution}{Unfortunately, full secret-independent execution is hardly achievable in practice. In fact, secret-independent execution does not only depends on the original implementation of an algorithm, but also on the multiple layers --- roughly 2: compilation/interpretation (\Cref{sec:programming_languages_and_software}) and execution (recall the concept of speculative execution and the Spectre and Meltdown attacks in \Cref{sec:security_of_cryptographic_algorithms:attacks:sca}) --- between the original implementation written by the developers and the corresponding instructions ultimately executed by the processor. As a corollary, it is better (although more costly) to check the secret-independent execution of compiled machine code rather than the original implementation.}

\subsubsection{Constant-Time Coding}\label{sec:security_of_cryptographic_algorithms:best_practices:constant}

\textbf{Constant-time coding} refers to writing code explicitly aimed at thwarting timing (and, collaterally, other side-channel) attacks. Some operations are naturally constant-time: many stream ciphers (\Cref{sec:cryptography_basics:ciphers}) --- such as ChaCha20 (\Cref{sec:symmetric_cryptography:ciphers:chacha20}) --- employ only bitwise operations like \texttt{xor}-ing, addition, and rotation; modern processors already execute these operations in constant time. 

\warningBlock{Misunderstanding the meaning of constant-time coding}{Constant-time coding \textbf{does not} prescribe that algorithms must execute in a fixed amount of time, but rather that secret data within an algorithm must not influence the use of underlying hardware resources.}

The inherent constant-time nature of bitwise operations is also exploited in the \textbf{bitslicing} technique \cite{mercadier_usuba_2018}, which consists of expressing a function in terms of single-bit operations (i.e., AND, XOR, OR, NOT) in a logic circuit and exploiting the bigger registries of modern processors (usually 64 bits) to parallelize the execution of such single-bit operations. For instance, an 8-bit value in a bitsliced implementation is stored in 8 variables (one for each bit), but a 64-bit architecture would allow processing up to 64-bit values at the same time.

To make an operation constant-time in presence of conditions (i.e., branches), it is common to compute both (or all) branches of the operation and then discard the useless result(s).\footnote{Go Assembly Mutation Testing (\url{https://words.filippo.io/assembly-mutation/})} An instance of this approach is the \textbf{Montgomery ladder} \cite{montgomery_speeding_1987} for scalar multiplication of a point $Q$ by a number $d$ in \gls{ecc} (\Cref{sec:mathematical_underpinnings:elliptic_curves,sec:trapdoor_functions:elliptic_curve_cryptography}). Simplifying, in the case of \gls{ecc}, the Montgomery ladder carries out the scalar multiplication by processing $d$ (represented in binary) one bit at a time. In general, point addition is done only if the bit is 1; however, the Montgomery ladder always performs point addition on a temporary variable (even if the bit is 0) and then checks through branching (e.g., an \texttt{if} conditional instruction) whether to disregard the temporary variable or swap its value with the actual variable. Notably, even such branching and swap need to be constant-time; this is done via \textbf{constant-time conditional swapping} usually implemented with bitwise XOR and AND bitwise operations.

\warningBlock{Consider other non-constant time operations}{Besides mathematical operations --- and, in general, code written directly by the developers of an algorithm --- developers should also check whether third-party code (e.g., libraries) is constant-time. For instance, comparisons --- e.g., checking equality among two sequences of bytes --- are usually not constant time but instead \textbf{fast failing}: they compare pairs of bytes until they differ, possibly not reaching the end of the two sequences. Hence, most implementations of comparisons leak which is the first byte where two sequences differ. Notably, some libraries are aware of the importance of constant-time coding and thereby offer constant-time functions (like the \texttt{CRYPTO\_memcmp} comparison function in OpenSSL --- see \url{https://docs.openssl.org/3.3/man3/CRYPTO_memcmp/}).}

\readBlock{More information on constant-time coding}{The Chosen Plaintext Consulting's and Red Hat's blog posts ``A beginner's guide to constant-time cryptography'' and ``The need for constant-time cryptography'' --- available at the links \url{https://www.chosenplaintext.ca/articles/beginners-guide-constant-time-cryptography.html} and \url{https://research.redhat.com/blog/article/the-need-for-constant-time-cryptography/} --- provide more information on (the relevance of) constant-time coding. Moreover, Thomas Pornin gives a good primer on constant-time cryptography in relation to his library BearSSL.\footnote{BearSSL - Constant-Time Crypto (\url{https://www.bearssl.org/constanttime.html})}}

\readBlock{A collection of constant-timeness verification tools}{The Centre for Research on Cryptography and Security at at Masaryk University curates a list of tools --- available at the link \url{https://crocs-muni.github.io/ct-tools/} --- for testing and verification of constant-timeness of programs.}

\subsubsection{Masking and Randomization}\label{sec:security_of_cryptographic_algorithms:best_practices:mask}

One of the simplest defenses against many side-channel attacks --- like timing, electromagnetic, and power analysis --- is to lower the \gls{snr} (\Cref{sec:security_of_cryptographic_algorithms:attacks:sca}) by \textbf{injecting noise}. For example, it is possible to execute operations in parallel on multiple cores while simultaneously adding dummy or decoy operations. However, lowering the \gls{snr} simply makes attacks harder but does not prevent them --- especially if measurements are precise and repeated enough times to filter out noise through statistics (i.e., differential attacks). Hence, noise injection is not a very effective defense.\footnote{\url{https://crypto.stackexchange.com/a/59901}}

Better (but not perfect) protection can be achieved through \textbf{masking}, which consists in mixing secret data with random values (e.g., through \texttt{xor}-ing) and using the masked secret data to execute operations --- intuitively, the operations have to be modified so that correct results can still be computed, and masks should never be reused. 

\exampleBlock{Examples of masking and randomization}{In \gls{rsa} it is possible to blind the message by multiplying it with a random number or the secret exponent by exploiting Fermat's little theorem.\footnote{\url{https://en.wikipedia.org/wiki/Fermat\%27s_little_theorem}} Similarly, in \gls{ecc} it is possible to exploit Jacobian and projective coordinates to blind point coordinates \cite{tunstall_coordinate_2010}.\footnote{\url{https://en.wikipedia.org/wiki/Jacobian_curve}}}

\subsubsection{Hardware Hardening}\label{sec:security_of_cryptographic_algorithms:best_practices:hardware}

Despite the fact that hardware-based cryptographic modules (\Cref{sec:cryptographic_modules}) are usually not enough to prevent (or sometimes even mitigate) the attacks discussed in \Cref{sec:security_of_cryptographic_algorithms:attacks}, their use is often considered a security best practice. In fact, even though being typically isolated from the main processors --- which reduces noise and gives the attacker cleaner and precise measurements --- \textbf{hardware-based cryptographic modules allows for implementing defenses not easily achievable through software alone}. For instance, many hardware-based cryptographic modules are shielded to reduce electromagnetic and sound leakage and embed power filtering to reduce power signal leakage. Moreover, randomizers for the power supply, the clock (called \textit{clock jitter}), and memory access (e.g., see Kilopass' Secretcode memory\footnote{\url{https://www.chipestimate.com/SIDE-CHANNEL-ATTACKS--How-Differential-Power-Analysis-DPA-and-Simple-Power-Analysis-SPA-Works-/Kilopass-Technology-a-part-of-Synopsys/Technical-Article/2017/09/19}}) are often available. Similarly, hardware-based cryptographic modules allows for balanced circuitry, that is, circuits that balance the amount of power used for a given data value or operation to prevent power analysis attacks. An even more sophisticated technique --- that comes at the cost of sensibly increased computational cost --- is \textbf{threshold implementation} and operates on a gate level: this technique exploits algorithms originally designed for \gls{smpc} (\Cref{sec:advanced_cryptography:secure_multi_party_computation}) to split secret data into multiple variables so that secret data are never used directly. Then, gates performing non-linear operations are replaced by a cluster of gates called \textbf{gadgets} which perform the same operation but have all inputs masked --- not dissimilarly to masking (\Cref{sec:security_of_cryptographic_algorithms:best_practices:mask}). Intuitively, threshold implementation is often not viable in practice.

\subsubsection{Memory Zeroization}\label{sec:security_of_cryptographic_algorithms:best_practices:zeroization}

To protect against data remanence attacks (\Cref{sec:security_of_cryptographic_algorithms:attacks:dra}) and memory leaks in general, it is recommended to immediately delete secret data when no longer needed. A common approach is to overwrite variables containing secret data with a string of zero bytes of the same length --- hence, \textbf{zeroization}. Another related recommendation is to minimize the time secret data is kept in memory and similarly minimize opportunities for dispersion. 

\remarkBlock{On the difficulty of achieving memory zeroization}{Unfortunately, memory zeroization is hindered by many factors. First, the programming language itself may not fully support zeroization (e.g., those based on garbage collection --- see \Cref{sec:programming_languages_and_software}). Also, processors may copy the value of variables into registries, stacks, or even swap them to non-volatile memory supports; it is rather difficult --- when not impossible --- to zeroize every copy of some secret data. Finally, processor optimizations may prevent zeroization altogether.}

Historically, several vulnerabilities found in the implementation of cryptographic algorithms were related to memory management in general (e.g., buffer overflows) which depends on the chosen programming language --- see \Cref{sec:programming_languages_and_software:languages}).


\subsubsection{Resistance to Fault Injection}\label{sec:security_of_cryptographic_algorithms:best_practices:fault}

A first strategy against fault injection attacks (\Cref{sec:security_of_cryptographic_algorithms:attacks:fia}) is to structure the code such that it \textbf{fails closed}: the execution should fall back to a closed (i.e., secure) state in all plausible fault scenarios, and only enter an open (i.e., insecure) state when an unlikely set of conditions are satisfied. In other words, it must never be possible to enter an insecure state by skipping or glitching any single instruction. A second strategy is \textbf{resiliency through redundancy}: critical operations should be performed more than once and their results compared, and the execution should safely fail in case inconsistencies are found. 

\exampleBlock{An example of resiliency through redundancy}{Boolean values may be replaced with multi-bit values with alternate bits set and used for checks.}
	
A third strategy is to introduce \textbf{random delays after externally observable events} (i.e., inter-chip communications, power-ups) to make it more difficult to time glitches.

\subsubsection{Key Limit}\label{sec:security_of_cryptographic_algorithms:best_practices:limit}

The probably most effective mitigation against side-channel attacks is to implement a hard limit on \textbf{key usage} --- a concept different than that of cryptoperiod (\Cref{sec:key_management:lifecycle:operational}). The intuition is that, by stopping using a key and replacing it before enough measurements about its usage can be gathered, it is possible to prevent the attacker from performing an effective analysis and recovering useful information. In fact, (especially differential) side-channel attacks require a statistically significant number of data points to be successful. Intuitively, key limit presents a trade-off between security and usability, as it forces frequent key rotations.

\subsubsection{Testing}\label{sec:security_of_cryptographic_algorithms:best_practices:test}

Any implementation --- especially that of cryptographic algorithms --- should be properly reviewed and tested for (both functional correctness and) security with one or (better) more among the following approaches:

\readBlock{More information on software security}{\Cref{sec:security_of_cryptographic_algorithms:best_practices:test} focuses on testing for the implementation of cryptographic algorithms; on the subject of generic software security, we refer the interested reader to the resources of the \gls{nist} \cite{scarfone_technical_2008,souppaya_secure_2022}, the \gls{iso} \cite{iso_24772_2024,iso_27034_2011}, the \gls{owasp} \cite{owasp_asvs_2024,owasp_wstg_2024}, and the \gls{enisa} \cite{enisa_advancing_2020}.}

\begin{itemize}
		
	\item \textbf{static analysis}: analyze the implementation without execution to detect vulnerabilities and ensure adherence to security best practices. While static analysis can detect vulnerabilities early in the development process and (ideally) achieve 100\% coverage, such an approach may produce false positives and relies on the expertise of the developers and the capabilities of the tools employed. As for examples of methodologies and tools for static analysis, KLEE\footnote{Documentation · KLEE (\url{https://klee-se.org/docs/})} \cite{cadar_klee_2008,erlingsson_efficient_2011} allows for checking memory safety and proper deallocation through symbolic execution, CodeSonar\footnote{Static Application Security Testing (SAST) Software Tool | CodeSonar (\url{https://codesecure.com/our-products/codesonar/})} provides thorough static analysis in source and binary code, and True Code\footnote{True Code - Automated embedded Software Security checks (\url{https://www.riscure.com/security-tools/true-code})} detects fault injection vulnerabilities;
	
	\item \textbf{dynamic analysis (attack simulation)}: execute the code and/or emulate real-world attacks to analyze the resulting runtime behavior under controlled conditions and identify vulnerabilities. In the latter case, developers assume the role of the attacker, with the noticeable advantage of already knowing the attack surface and operating in favorable conditions (e.g., not needing depackaging/decapsulation --- see \Cref{sec:security_of_cryptographic_algorithms:attacks:sca}). While attack simulation provides practical evidence of eventual vulnerabilities, the coverage of such an approach is limited by the developers' offensive capabilities. As for examples of methodologies and tools for dynamic analysis, project Wycheproof\footnote{Project Wycheproof \url{https://github.com/C2SP/wycheproof}} contains test vectors against known attacks, ctgrind\footnote{ctgrind (\url{https://github.com/agl/ctgrind?tab=readme-ov-file})} checks whether code is constant-time, Inspector SCA,\footnote{Side Channel Analysis Security Solutions - Riscure (\url{https://www.riscure.com/security-tools/inspector-sca/})} Inspector FI,\footnote{Fault Injection Security Tools - Riscure (\url{https://www.riscure.com/security-tools/inspector-fi/})}, and ChipWhisperer\footnote{ChipWhisperer - the complete open-source toolchain for side-channel power analysis and glitching attacks (\url{https://github.com/newaetech/chipwhisperer})} allows for conducting side-channel and fault injection attacks;
		
	\item \textbf{formal methods}: mathematically prove the security of the implementation against one or more attacks. Formal methods usually require modeling the system under test and specifying the desired security properties in a formal language, using a tool for conducting automated analysis and identifying (proofs of) vulnerabilities, and iteratively refinement. While formal methods offer mathematical assurance of security, such an approach can be resource-intensive and require deep expertise on the developers' end. As for examples of methodologies and tools for formal methods, MaskVerif\footnote{maskVerif (\url{https://cryptoexperts.com/maskverif/})} \cite{sako_maskverif_2019} provides automatic verification of side-channel countermeasures based on masking, EasyCrypt\footnote{EasyCrypt (\url{https://www.easycrypt.info/})} \cite{hutchison_computer_2011,aldini_easycrypt_2014} manages security proofs of cryptographic constructions, Vale\footnote{Introduction — Vale: Verified Assembly Language for Everest documentation (\url{https://project-everest.github.io/vale/intro.html})} allows for writing in a C-like formal pseudo-language and then directly generate assembly code, ct-verif\footnote{GitHub - imdea-software/verifying-constant-time (\url{https://github.com/imdea-software/verifying-constant-time})} \cite{almeida_verifying_2016} validates constant-time coding, and FIVER\footnote{GitHub - Chair-for-Security-Engineering/FIVER (\url{https://github.com/Chair-for-Security-Engineering/FIVER})} \cite{richter_fiver_2021} analyzes and verifies implementations of countermeasures against fault-injection attacks;
	\readBlock{More information on formal methods for cryptography}{A survey of approaches and tools used by cryptographers and their features can be found in \cite{barbosa_sok_2021}.}
	
	\item \textbf{fuzzing}: automatically generate and test a large number of inputs to identify unexpected behavior or vulnerabilities in the implementation. While fuzzing is inherently scalable and can uncover edge cases, such an approach may also not detect all security issues, especially those requiring specific attack vectors. As for examples of methodologies and tools for fuzzing, AFL (American Fuzzy Lop)\footnote{GitHub - google/AFL: american fuzzy lop - a security-oriented fuzzer (\url{https://github.com/google/AFL})} allows for fuzzing generic software, LibFuzzer\footnote{libFuzzer – a library for coverage-guided fuzz testing. — LLVM 20.0.0git documentation (\url{https://llvm.org/docs/LibFuzzer.html})} provides in-process, coverage-guided, evolutionary fuzzing, FuzzTest \footnote{google/fuzztest (\url{https://github.com/google/fuzztest})} allows for writing and executing fuzz tests for C++ software, and True Code simulates faults in virtual hardware;
		
	\item \textbf{machine learning and artificial intelligence}: detect vulnerabilities (e.g., by automatically analyzing logs, traces, and execution flows for suspicious patterns), optimize attack strategies (e.g., by suggesting likely weak spots), and analyze the behavior of the implementation to identify anomalies. While machine learning and artificial intelligence can uncover non-obvious vulnerabilities and minimize the developers' effort, such an approach requires high-quality data for training and significant computational resources. SCAAML (Side Channel Attacks Assisted with Machine Learning)\footnote{GitHub - google/scaaml (\url{https://github.com/google/scaaml/})} is an example of a tool based on machine learning and artificial intelligence;
		
	\item \textbf{compliance testing}, which aims to validate the security of the implementation against established security standards (\Cref{sec:standards_and_certifications}).
		
\end{itemize}

Developers should carefully assess the most suitable approach according to their context (e.g., expertise, threat model --- see \Cref{sec:cryptographic_security_overview:threat_model} --- and programming language --- see \Cref{sec:programming_languages_and_software}) and possibly combine multiple approaches. For instance, developers may use static analysis and formal methods for verifying foundational security properties and attack simulation and fuzzing for practical testing of the implementation --- together with machine learning and artificial intelligence to enhance attack simulations and analysis, and compliance testing for ensuring adherence to industry standards, regulatory requirements, and organizational security policies. 

\readBlock{More information on testing for cryptographic code}{Further resources relevant to the topic of testing the security of the implementation of cryptographic algorithms can be found in \cite{bhargavan_everest_2017,buhan_sok_2022,piscitelli_fault_2015,mccann_towards_2017,veshchikov_use_2017}.}

\readBlock{Evaluating side-channel and fault injection attacks resistance}{The German \gls{bsi} has published the AIS 46 series of guidelines --- available at the link \url{https://www.bsi.bund.de/EN/Themen/Unternehmen-und-Organisationen/Informationen-und-Empfehlungen/Kryptografie/Seitenkanalresistenz/seitenkanalresistenz_node.html} --- on evaluating side-channel and fault injection resistance of algorithm implementation.}

\subsubsection{Bad Practices}\label{sec:security_of_cryptographic_algorithms:best_practices:bad:practices}

Besides the aforementioned best practices, there exist also bad practices in implementing algorithms. The most notorious are:

\begin{itemize}
	
	\item \textbf{security through obscurity}, already discussed in \Cref{sec:cryptography_basics};
		
	\item \textbf{rolling your own crypto}, that is, programmers who prefer developing their own implementation of algorithms rather than relying on well-established and vetted cryptographic libraries. After reading \Cref{sec:security_of_cryptographic_algorithms}, it should be clear how difficult it is to implement algorithms securely --- not to mention bugs and inefficiencies. For instance, an empirical study of vulnerabilities in cryptographic libraries (``You Really Shouldn't Roll Your Own Crypto'') found that only 27.2\% of vulnerabilities in cryptographic libraries are actually cryptographic-related issues, while systems-level bugs are a greater security concern with memory safety issues (\Cref{sec:programming_languages_and_software:languages}) constituting 37.2\% of vulnerabilities \cite{blessing_you_2021}. In fact, note that there are experts who pursue security in algorithm implementation as a full-time profession;\footnote{\url{https://words.filippo.io/full-time-maintainer/}}
	
	\item \textbf{bloated code base}, that is: the more the code, the more vulnerabilities and errors are present. For this reason, modern cryptographic libraries tend to have a small code base (e.g., TweetNaCl\footnote{\url{https://tweetnacl.cr.yp.to/}} has ~1,000 lines of code) hence allowing for complete and thorough security audits. Also, a larger code base usually implies many \glspl{api}, increasing the risk of misusing and misconfiguration --- especially by non-expert developers.

\end{itemize}

\readBlock{Further low-level bad practices}{Greg Rubin's blog post ``Crypto Gotchas!'' --- available at the link \url{https://gotchas.salusa.dev/} --- provides a very interesting list of low-level bad practices concerning applied cryptography.}
